% k-tunable ablation — power-of-choices family (Po2 default, Po4, Po8, Fanout)
\subsection{\texorpdfstring{$k$-tunable Dispatch: from Po2 to Fanout}{k-tunable Dispatch: from Po2 to Fanout}}\label{section:po2_ablation}

Astrolabe's core dispatch policy is \emph{power-of-$k$ choices}: for each request, sample $k$ random instances, query each Predictor in parallel, and route to the one with the lowest predicted latency. Increasing $k$ trades higher control-plane fanout cost (linear in $k$) for better load-balancing quality and stronger mitigation of herding~\cite{mitzenmacher2001,10.1137/S0097539795288490}. We evaluate $k \in \{2, 4, 8, 12\}$; $k{=}12$ recovers full Fanout.

\begin{figure}[!tb]
    \centering
    \includegraphics[width=0.90\linewidth]{figure/revision/po2_comparison.png}
    \caption{Power-of-$k$ choices ablation}
    \Description{Capacity, E2E latency, throughput, and scheduling overhead for power-of-k choices variants from k=2 to full fanout.}
    \label{fig:po2_comparison}
\end{figure}

Figure~\ref{fig:po2_comparison} shows capacity, mean/P99 E2E latency at each scheduler's own capacity, throughput, and per-request scheduling overhead. Capacity is effectively tied from $k{=}2$ upward: Po2-est reaches \capPoTwoEst{} QPS vs Po4-est \capPoFourEst{} (the maximum), Po8-est \capPoEightEst{}, and Fanout-est \capFanoutEst{}, all within a $\pm$0.3\,QPS band at SLO\,=\,\sloSec{}s. Scheduling overhead grows with $k$ through $k{=}8$ and then plateaus, because each dispatch fans out to more predictors and is bounded by the slowest reply.

Randomized probing also mitigates herding. When many bursty requests arrive within one round trip, a deterministic Fanout scheduler can compute identical predictions and overload the same instance, while a power-of-$k$ policy breaks this correlation by randomizing each request's candidate set. This helps the sampled variants maintain lower E2E latency and higher throughput than Fanout in Figure~\ref{fig:po2_comparison}. In this sweep, $k{=}4$ gives the best empirical tradeoff: it has the highest capacity and throughput, while retaining the second-lowest E2E latency and scheduling overhead. We nevertheless use Po2 as the default in the main evaluation to follow the standard power-of-two-choices convention without tuning $k$ on the evaluation workload.
